\documentclass[12pt,reqno]{amsart}
\usepackage{../handout, ../customcommands}
\usepackage{caption}

\hdr{MAT 240}{\S3.7 Multivariate distributions, part 4}

\begin{document}

\textbf{Suggested problems:} 1, 3, 6, 8

\bigskip
\hrule

\bigskip
\prob Define the function $\sigma: \mathbb{R} \to \mathbb{R}$ by setting
\[
	\sigma(z) = \frac{1}{1+e^{-z}}.
\]
This is called the \textit{sigmoid function}. 

\medskip
\begin{enumerate}
\item Draw a graph of $\sigma$.\vfill
\item Compute the range of $\sigma$.\vfill
\end{enumerate}

\prob Suppose that $X_1$ and $X_2$ are independent indicator random variables, both distributed as $\mathcal{B}er(1/2)$, where $X_1$ codes whether a patient smokes and $X_2$ codes whether they have family history of a certain disease. Suppose $Y$ is a third indicator random variable that codes whether the patient has the disease. Suppose that we model $Y$ with a logistic regression model with predictors $X_1$ and $X_2$, and parameters $\beta_0 = -2$, $\beta_1 = 1$, and $\beta_2 = 1.5$.

\begin{enumerate}
\item Compute the marginal p.m.f.\ $f(y)$.\vfill
\item Given that a patient has the disease, what is the probability that they both smoke and have a family history?\vfill
\item Given that a patient has the disease, what is the probability that they smoke?\vfill
\end{enumerate}


\end{document}